% Copyright (c) 2026 Eugene Joseph M. Ragasa.
% SPDX-License-Identifier: Apache-2.0
% Author: Eugene Joseph M. Ragasa
% Formal proof completed: 2026-10-30
\documentclass[11pt]{article}
\usepackage[T1]{fontenc}
\usepackage{amsmath,amssymb}
\usepackage{booktabs}
\usepackage[margin=1in]{geometry}
\usepackage[hidelinks]{hyperref}
\usepackage{tabularx}

\title{Proof Dossier: Reconstruction After Negative-Nyquist Deletion}
\author{Eugene Joseph M. Ragasa}
\date{Formal proof completed: 2026-10-30}

\begin{document}
\raggedright
\maketitle

\section{Identity and status}
\begin{tabularx}{\textwidth}{@{}lX@{}}
\toprule
Statement ID & \texttt{BG-FOURIER-THM-086} \\
Status & proved; no admissions \\
Lean declaration & \path{ProjectKoios.BandGauge.Fourier.finiteReconstruction_eq_negativeNyquistRemoved_add} \\
Owning module & \path{ProjectKoios/BandGauge/Fourier/CenteredEvenMesh.lean} \\
Author and license & Eugene Joseph M. Ragasa; \texttt{Apache-2.0} \\
\bottomrule
\end{tabularx}

\section{Exact mathematical statement}
Let $M$ be a positive even centered mesh with finite mode type $\mathcal M$ and distinguished
negative-Nyquist mode $\nu_-$.  Let $V$ be an additive commutative complex module, let
$\phi:Q\times\mathcal M\to\mathbb C$ be any synthesis kernel, and let
$c:\mathcal M\to V$ be any coefficient field.  For every $q\in Q$,
\[
 \sum_{m\in\mathcal M}\phi(q,m)c_m
 =\sum_{m\in\mathcal M\setminus\{\nu_-\}}\phi(q,m)c_m
   +\phi(q,\nu_-)c_{\nu_-}.
\]
Scalar multiplication is understood in $V$.  The Lean declaration states this with
\texttt{finiteReconstruction} and \texttt{negativeNyquistRemovedReconstruction}; it assumes only
\texttt{AddCommMonoid V} and \texttt{Module Complex V}.

\section{Immutable source provenance}
The exact declaration was introduced at Formal commit
\texttt{458e480e08a925d373f352f76c0dd4fe1fc356d7}, reviewed and merged through PR~\#26 at
\texttt{ad5fd4472d3bbdabe63a976dcf148f8a3eaee2d5}.  The owning module SHA-256 is
\texttt{ce13025abea15f70a6754cccf9568984150223d90c28bb62c071875fae223aca}; the exact locator is
\texttt{BG-FOURIER-THM-086}.

Motivation came from the centered-even convention inspected read-only at KSD PR~\#80 merge
\texttt{d8e7d1db3eaf16332e691427fb9d5949325badd6}, tree
\texttt{23e6b97f3eddb074ab190f2576a4c2fff24e29f6}.  The owner paths, hashes, and dirty-checkout
limitation are in \path{docs/source-audit.md}.  The relation is motivational only: the theorem is
kernel-agnostic and does not bind KSD execution or bytes.

\section{Assumptions, conventions, and statement delta}
\begin{tabularx}{\textwidth}{@{}lXX@{}}
\toprule
Item & Lean encoding & Boundary \\
\midrule
Mode inventory & finite centered-mode subtype & exact abstract type, not a runtime array \\
Removed support & full inventory minus $\nu_-$ & one deletion; no coefficient redistribution \\
Kernel & arbitrary complex-valued function $\phi$ & no Fourier character or normalization assumption \\
Coefficient carrier & additive commutative complex module $V$ & no norm, topology, matrix rank, or Hermiticity needed \\
Sign & full equals removed plus deleted term & not a representation-preserving equivalence \\
\bottomrule
\end{tabularx}
The result is stronger in algebraic generality but deliberately weaker in application meaning than
a statement about a concrete FFT reconstruction.

\section{Statement-specific prior art and contribution boundary}
Finite-sum partition by a retained support and its complement is standard and
\texttt{background-only}.  No novelty or priority claim is made for the algebra.  The Project Koios
contribution is to expose the exact missing term rather than silently treating negative-Nyquist
deletion as an equivalent centered representation.

\section{Proof structure and mathematical argument}
The generic retained--omitted decomposition gives
\[
 \sum_{m\in\mathcal M}\phi(q,m)c_m
 =\sum_{m\in S_-}\phi(q,m)c_m
  +\sum_{m\in S_-^{\mathsf c}}\phi(q,m)c_m.
\]
By \texttt{BG-FOURIER-THM-081}, $S_-^{\mathsf c}=\{\nu_-\}$.  A sum over this singleton is
$\phi(q,\nu_-)c_{\nu_-}$.  Substituting the definition of the removed reconstruction produces the
stated identity.  No limiting or norm argument occurs.

\section{Lean realization and dependencies}
Lean first applies
\path{finiteReconstruction_eq_truncated_add_omitted}.  Simplification then unfolds
\texttt{negativeNyquistRemovedReconstruction}, rewrites the omitted complement using the singleton
result, and evaluates the singleton sum.  The critical local dependency is
\texttt{BG-FOURIER-THM-081}; the generic decomposition is in
\path{ProjectKoios/BandGauge/Fourier/FiniteReconstruction.lean}.

Toolchain: Lean \texttt{4.35.0}, commit
\texttt{eb6cafb7290e8bd5f2937522d8d804f58f02bafc}; Mathlib
\texttt{56a005855ea17e0b0958f73ec801ae140b0a1679}.

\section{Trusted base and axiom report}
\begin{verbatim}
#print axioms
  ProjectKoios.BandGauge.Fourier.
    finiteReconstruction_eq_negativeNyquistRemoved_add
-- [propext, Classical.choice, Quot.sound]
\end{verbatim}
The repository admission audit reports no repository-owned \texttt{axiom}, \texttt{sorry}, or
\texttt{admit} for this declaration.

\section{AI, automation, and review record}
The Lean proof and exposition were AI-assisted and checked against the source and theorem manifest
by the maintainer.  Lean kernel checking, simplification, build output, hashing, and LaTeX
compilation are automation, not independent expert review.  No independent mathematical,
scientific, novelty, or publication review is recorded.

\section{Scientific, numerical, and certificate boundaries}
This exact identity is not an off-node reconstruction bound and supplies no maximum or
root-mean-square error.  It does not show that binary64 runtime coefficients equal $c$, that a
runtime phase equals $\phi$, or that a mesh represents a physical reciprocal path.  It proves no
ordinary omitted-tail claim, spectral adequacy, material applicability, or checked certificate.

\section{Unresolved questions}
A future exact/outward owner bridge must separately authenticate the runtime inventory, coefficient
field, kernel, sign, phase, and deletion operation.  Any split or redistributed Nyquist coefficient
would require a new reconstruction theorem rather than this singleton-deletion identity.

\section{Review and revision history}
\begin{tabularx}{\textwidth}{@{}lllX@{}}
\toprule
Date & Revision & Process & Outcome and limitation \\
\midrule
2026-10-30 & initial dossier & AI-assisted maintainer review & exact decomposition and boundaries checked; no independent expert review \\
\bottomrule
\end{tabularx}

\end{document}
